% =====================================================================
%  STA257 — Probability and Statistics I
%  Note-Taking Template, organized by Chapter / Section (1, 1.1, 1.2 ...)
%  --------------------------------------------------------------------
%  Soft, muted accent colors (not saturated blocks): Definition (blue),
%  Theorem (amber/yellow), Example (green). Everything else (Lemma,
%  Proposition, Corollary, Remark, Proof) is PLAIN TEXT with a quiet
%  left rule — bold title only, no box / background.
%
%  Numbering: \section = 1, 2, 3 ...   \subsection = 1.1, 1.2 ...
%  Boxes are numbered within the section: Definition 1.1, Theorem 1.2 ...
% =====================================================================

\documentclass[11pt]{article}

% ---------- Page & Fonts ----------
\usepackage[margin=1in]{geometry}
\usepackage[utf8]{inputenc}
\usepackage[T1]{fontenc}
% \usepackage{lmodern}      % enable on Overleaf / full TeX installs
% \usepackage{mathpazo}     % nicer serif math font; enable if available
% \usepackage{microtype}    % enable if available — nicer justification

% ---------- Math ----------
\usepackage{amsmath, amssymb, amsfonts}
\usepackage{amsthm}
\usepackage{mathtools}
\usepackage{bm}

% ---------- Colors & Boxes ----------
\usepackage{xcolor}
\usepackage[most]{tcolorbox}

% ---------- Misc ----------
\usepackage{enumitem}
\usepackage{titlesec}
\usepackage{tikz}
\usepackage{framed}
\usepackage{etoolbox}
\usepackage{hyperref}
\usetikzlibrary{arrows.meta}

% ---------- Color palette (soft / muted) ----------
\definecolor{defColor}{HTML}{5B8DEF}    % soft blue   - Definition
\definecolor{thmColor}{HTML}{E0A030}    % soft amber  - Theorem
\definecolor{exColor}{HTML}{4CA37A}     % soft green  - Example
\definecolor{neutralColor}{HTML}{8B8F98}% soft gray   - text accents
\definecolor{inkColor}{HTML}{2D3138}    % near-black text

% ---------- Section numbering: 1, 1.1, 1.2, 2, 2.1 ... ----------
\setcounter{secnumdepth}{2}
\setcounter{tocdepth}{2}

% =====================================================================
% ---------- The 3 "headline" boxes — soft tinted style ----------
%  All boxes share ONE counter, numbered within \section:
%  Definition 1.1, Theorem 1.2, Example 1.3, Lemma 1.4 ...
% =====================================================================
\newtcbtheorem[number within=section]{defbox}{Definition}%
{colback=defColor!5, colframe=defColor!65, fonttitle=\bfseries,
 coltitle=defColor!75!black, colbacktitle=defColor!12,
 boxrule=0.6pt, arc=1.5mm,
 left=6pt, right=6pt, top=4pt, bottom=4pt,
 breakable, enhanced jigsaw}{def}

\newtcbtheorem[use counter from=defbox]{thmbox}{Theorem}%
{colback=thmColor!6, colframe=thmColor!70, fonttitle=\bfseries,
 coltitle=thmColor!60!black, colbacktitle=thmColor!15,
 boxrule=0.6pt, arc=1.5mm,
 left=6pt, right=6pt, top=4pt, bottom=4pt,
 breakable, enhanced jigsaw}{thm}

\newtcbtheorem[use counter from=defbox]{exbox}{Example}%
{colback=exColor!5, colframe=exColor!65, fonttitle=\bfseries,
 coltitle=exColor!70!black, colbacktitle=exColor!12,
 boxrule=0.6pt, arc=1.5mm,
 left=6pt, right=6pt, top=4pt, bottom=4pt,
 breakable, enhanced jigsaw}{ex}

% 只要一个 {标题}，标签自动留空；标题可以写 {} 留空
\newenvironment{definition}[1]{\begin{defbox}{#1}{}}{\end{defbox}}
\newenvironment{theorem}[1]{\begin{thmbox}{#1}{}}{\end{thmbox}}
\newenvironment{example}[1]{\begin{exbox}{#1}{}}{\end{exbox}}

% =====================================================================
% ---------- Everything else: PLAIN (no fill, no frame), with a quiet
%  vertical rule + generous spacing. Shares the same counter as boxes.
% =====================================================================
\newtheoremstyle{quiet}%
  {12pt}{12pt}%          space above / below
  {\normalfont}{0pt}%    body font / indent
  {\bfseries}{.}%        head font / punctuation after head
  {0.6em}{}%             space after head / custom head-spec
\theoremstyle{quiet}

\makeatletter
\newtheorem{proposition}[tcb@cnt@defbox]{Proposition}
\newtheorem{lemma}[tcb@cnt@defbox]{Lemma}
\newtheorem{corollary}[tcb@cnt@defbox]{Corollary}
\makeatother

\newtheoremstyle{quietremark}%
  {12pt}{12pt}%
  {\normalfont\itshape}{0pt}%
  {\bfseries\itshape}{.}%
  {0.6em}{}
\theoremstyle{quietremark}
\newtheorem*{remark}{Remark}

% -- the quiet left-rule "bar" --
\definecolor{barColor}{HTML}{8B8F98}
\newenvironment{quietbar}{%
  \def\FrameCommand{{\color{barColor!55}\vrule width 1.1pt\hspace{9pt}}}%
  \MakeFramed{\advance\hsize-\width\relax \FrameRestore}%
}{%
  \endMakeFramed
}
% Auto-wrap every plain statement (and proofs) in the quiet bar.
\BeforeBeginEnvironment{lemma}{\begin{quietbar}}
\AfterEndEnvironment{lemma}{\end{quietbar}}
\BeforeBeginEnvironment{proposition}{\begin{quietbar}}
\AfterEndEnvironment{proposition}{\end{quietbar}}
\BeforeBeginEnvironment{corollary}{\begin{quietbar}}
\AfterEndEnvironment{corollary}{\end{quietbar}}
\BeforeBeginEnvironment{remark}{\begin{quietbar}}
\AfterEndEnvironment{remark}{\end{quietbar}}
\BeforeBeginEnvironment{proof}{\begin{quietbar}}
\AfterEndEnvironment{proof}{\end{quietbar}}

% Proof ends with a small solid square
\renewcommand{\qedsymbol}{$\blacksquare$}

% ---------- Section styling ----------
\titleformat{\section}
 {\normalfont\Large\bfseries\color{inkColor}}
 {\thesection}{0.6em}{}
 [{\color{neutralColor!35}\titlerule[0.8pt]}]
\titleformat{\subsection}
 {\normalfont\large\bfseries\color{inkColor}}
 {\thesubsection}{0.6em}{}
\titleformat{\subsubsection}
 {\normalfont\normalsize\bfseries\color{inkColor}}
 {\thesubsubsection}{0.6em}{}

% ---------- Hyperref colors ----------
\hypersetup{
 colorlinks=true,
 linkcolor=inkColor,
 citecolor=exColor,
 urlcolor=defColor,
 pdftitle={STA257 Notes},
 pdfauthor={Fuxuan Xu}
}

% ---------- Quick math macros (probability) ----------
\newcommand{\R}{\mathbb{R}}
\newcommand{\N}{\mathbb{N}}
\newcommand{\Z}{\mathbb{Z}}
\newcommand{\Q}{\mathbb{Q}}
\newcommand{\eps}{\varepsilon}
\newcommand{\abs}[1]{\left\lvert #1 \right\rvert}
\newcommand{\dd}{\mathrm{d}}

\renewcommand{\P}{\mathbb{P}}          % probability   \P(A)
\newcommand{\E}{\mathbb{E}}            % expectation   \E[X]
\DeclareMathOperator{\Var}{Var}        % variance      \Var(X)
\DeclareMathOperator{\Cov}{Cov}        % covariance    \Cov(X,Y)
\DeclareMathOperator{\Corr}{Corr}      % correlation   \Corr(X,Y)
\newcommand{\1}{\mathbf{1}}            % indicator     \1_A
\newcommand{\iid}{\overset{\text{i.i.d.}}{\sim}}
\newcommand{\cond}{\,\middle|\,}       % use inside \left( ... \right)
\newcommand{\toP}{\xrightarrow{\;P\;}}         % convergence in probability
\newcommand{\toas}{\xrightarrow{\;\text{a.s.}\;}}% almost sure convergence
\newcommand{\toD}{\xrightarrow{\;D\;}}         % convergence in distribution

% Common distributions:  X \sim \Bin(n,p),  \Normal(\mu,\sigma^2) ...
\DeclareMathOperator{\Bern}{Bernoulli}
\DeclareMathOperator{\Bin}{Binomial}
\DeclareMathOperator{\Geom}{Geometric}
\DeclareMathOperator{\NegBin}{NegBinomial}
\DeclareMathOperator{\Pois}{Poisson}
\DeclareMathOperator{\Hyper}{Hypergeometric}
\DeclareMathOperator{\Unif}{Uniform}
\DeclareMathOperator{\Expo}{Exponential}
\DeclareMathOperator{\GammaDist}{Gamma}
\DeclareMathOperator{\BetaDist}{Beta}
\newcommand{\Normal}{N}

% =====================================================================
% ---------------------------- COVER PAGE ------------------------------
% =====================================================================
\begin{document}

\begin{titlepage}
\centering
\vspace*{1.3cm}

{\Large\textsc{Lecture Notes}}\\[0.4cm]
{\color{neutralColor!50}\rule{0.35\textwidth}{0.6pt}}\\[1.0cm]

{\fontsize{32}{38}\selectfont\bfseries STA257}\\[0.3cm]
{\Large\itshape Probability and Statistics I}\\[1.0cm]

% ---------------- Probability illustration ----------------
\begin{tikzpicture}[scale=1.0]
  % axes
  \draw[neutralColor!70, -Latex] (-3.2,0) -- (3.4,0) node[right, font=\tiny, inkColor] {$x$};
  \draw[neutralColor!70, -Latex] (-3.0,-0.1) -- (-3.0,2.6);

  % shaded area under the bell curve:  P(a < X < b)
  \fill[defColor!18, domain=-0.8:1.3, smooth, variable=\x]
    (-0.8,0) -- plot ({\x},{2.2*exp(-\x*\x/1.2)}) -- (1.3,0) -- cycle;
  \draw[defColor!70, thick] (-0.8,0) -- (-0.8,{2.2*exp(-0.64/1.2)});
  \draw[defColor!70, thick] (1.3,0) -- (1.3,{2.2*exp(-1.69/1.2)});
  \node[defColor!70!black, font=\tiny] at (0.25,0.55) {$\P(a<X<b)$};
  \node[inkColor, font=\tiny] at (-0.8,-0.22) {$a$};
  \node[inkColor, font=\tiny] at (1.3,-0.22) {$b$};

  % the normal density curve
  \draw[neutralColor!80, thick, domain=-3:3, smooth, samples=80, variable=\x]
    plot ({\x},{2.2*exp(-\x*\x/1.2)});

  % mean marker
  \draw[thmColor!70, dashed, thick] (0,0) -- (0,2.2);
  \node[thmColor!55!black, font=\tiny] at (0.22,2.35) {$\mu$};
  \node[thmColor!55!black, font=\tiny] at (2.25,1.6) {$f_X(x)$};

  % --- arrow to a small discrete pmf (bar chart) ---
  \draw[exColor!70, -Latex, thick] (3.0,1.4) .. controls (3.7,1.9) and (4.3,1.6) .. (4.3,1.15);
  \node[exColor!60!black, font=\tiny] at (3.85,2.05) {discrete};
  \foreach \k/\h in {0/0.25, 1/0.7, 2/0.95, 3/0.6, 4/0.25}{
    \fill[exColor!25] ({3.75+0.28*\k},0) rectangle ({3.95+0.28*\k},\h);
    \draw[exColor!70] ({3.75+0.28*\k},0) rectangle ({3.95+0.28*\k},\h);
  }
  \draw[exColor!70] (3.65,0) -- (5.25,0);
  \node[exColor!60!black, font=\tiny] at (4.45,-0.25) {$p_X(k)$};
\end{tikzpicture}

\vspace{1.7cm}
{\large Fuxuan Xu}\\[0.2cm]
{\large University of Toronto}

\vfill

\begin{tcolorbox}[colback=black!2, colframe=neutralColor!35, width=0.7\textwidth,
 arc=1.5mm, boxrule=0.4pt]
\centering\small
\textbf{How this notebook works}\\[3pt]
Notes are organized \textbf{by topic} --- chapters (1, 2, \dots) and sections
(1.1, 1.2, \dots). Only three box types carry
color --- \textcolor{defColor!75!black}{\textbf{Definition}} (blue),
\textcolor{thmColor!60!black}{\textbf{Theorem}} (amber), and
\textcolor{exColor!70!black}{\textbf{Example}} (green) --- kept soft and light
so they're easy on the eyes. Everything else (Lemma, Proposition, Corollary,
Remark, Proof) is plain text with a bold label, no box.
\end{tcolorbox}

\vspace{1.0cm}
\end{titlepage}

\tableofcontents
\newpage

% =====================================================================
\section{Probability}
% =====================================================================

\subsection{Probability Space}

\subsection{The Empirical Law of Large Numbers}

\subsection{Uniform Distribution}

\subsection{Exponential Distribution}

\subsection{Probability Density Functions}

\subsection{Conditional Probability}
\begin{definition}{Conditional Probability}
For any two events, $A$ and $B$, for a random variable $W$ defined on a sample space $\Omega$, the conditional probability of $A$ given $B$ is:
\[
P(W\in A|W\in B)=\frac{P(W\in AB)}{P(W\in B)},
\]
which means that:
\[
P(A|B)=\frac{P(AB)}{P(B)},
\]
and that:
\[
P(AB)=P(A|B)P(B)=P(B|A)P(A).
\]
\end{definition}

\begin{theorem}{Bayes' Theorem}
Given a partition of the sample space $\Omega$ intro non-trivial mutually disjoint events $B_n$, where $n\in\mathbb{N}$, then for any other event $A\subset\Omega$, we obtain:
\begin{enumerate}
\item Total probability:
\[
\displaystyle P(A)=\sum_{n\in\mathbb{N}}P(A|B_n)P(B_n).
\]
\item Bayes' formula:
\[
\displaystyle P(B_n|A)=\frac{P(A|B_n)P(B_n)}{P(A)}.
\]
\end{enumerate}
\end{theorem}

\begin{proof}
\begin{enumerate}
\item Since $A=A\Omega=A\sum_{n=1}^\infty B_n$, we have that:
\[
P(A)=\sum_{n=1}^\infty P(AB_n)=\sum_{n=1}^\infty P(A|B_n)P(B_n).
\]
\item Since $P(AB_n)=P(A|B_n)P(B_n)$, we have that:
\[
\frac{P(AB_n)}{P(A)}=P(B_n|A)=\frac{P(A|B_n)P(B_n)}{P(A)}.
\]
\end{enumerate}
\end{proof}


\subsection{Monte Carlo Method}

\newpage
% =====================================================================
\section{Expectation}
% =====================================================================

\subsection{The Linearity of Expectation}

\subsection{Finite Discrete Distribution}

\subsection{The Bernoulli Trial}

\subsection{The Continuous Case}

\subsection{Density and Expectation}

\subsection{Conditional Expectation}

\newpage
% =====================================================================
\section{Geometry of Data}
% =====================================================================

\subsection{The Geometry of Random Variables}
\begin{definition}{Inner Product, Norm and Metric}
The three basic elements of Euclidean geometry, expressed in random variables and expectations:
\begin{enumerate}
\item Inner product: $\displaystyle\frac{\langle x,y\rangle}{n}=\frac{\sum_{i=1}^n x_iy_i}{n}=EXY$.
\item Norm: $\displaystyle\frac{\|x\|}{n}=\sqrt{\frac{\langle x,x\rangle}{n}}=\sqrt{EX^2}.$
\item Metric: $\displaystyle\frac{d(x,y)}{\sqrt{n}}=\frac{\|y-x\|}{\sqrt{n}}=\sqrt{E(Y-X)^2}$.
\end{enumerate}
\end{definition}

\begin{definition}{$L$ space}
The random variable space $L$ is a Banach space such that:
\[
L=\{X:E|X|<\infty\},
\]
where $L$ is a complete metric space, but the inner product does not exist in $L$.
\end{definition}

\begin{definition}{$L_2$ Space}
The random variable space $L_2$ is a Hilbert space, and a subspace of $L$, where:
\[
L_2=\{X\in L\mid EX^2<\infty\},
\]
and the inner product, norm and metric are well-defined on $L_2$.
\end{definition}

\begin{lemma}
Let $X,Y$ be random variables, if $X,Y\in L_2$, then $X,Y\in L$ and $XY\in L$.
\end{lemma}

\begin{proof}
Let $X,Y\in L_2$, then $EX^2<\infty,~EY^2<\infty$, therefore we have:
\[
\displaystyle E|X|<\frac{1+EX^2}{2}<\infty, E|Y|<\frac{1+EY^2}{2}<\infty.
\] 
Hence $X,Y\in L$.\\
Since $\displaystyle|XY|\le \frac{X^2+Y^2}{2}$, we have:
\[
\displaystyle E|XY|\le E\left(\frac{X^2+Y^2}{2}\right)=\frac{EX^2}{2}+\frac{EY^2}{2}<\infty.
\]
Hence, $XY\in L$.
\end{proof}

\begin{lemma}
Let $X,Y$ be random variables, if $X,Y\in L_2$, then $X+Y\in L_2$.
\end{lemma}

\begin{proof}
Let $X,Y\in L_2$, then $EX^2<\infty, EY^2<\infty$. By Cauchy-Schwarz:
\[
EXY\le E|XY|\le\sqrt{EX^2EY^2}<\infty.
\]
Hence, we have:
\[
\begin{aligned}
E(X+Y)^2&=E(X^2+2XY+Y^2)\\&=EX^2+2EXY+EY^2\\&<\infty.
\end{aligned}
\]
Hence, $X+Y\in L_2$.
\end{proof}

\begin{lemma}{Properties of $EXY$ as an Inner Product}
The expectation $EXY$ satisfies the following three properties of an inner product:
\begin{enumerate}
\item Bilinearity: $\langle a_1X_1+a_2X_2,Y\rangle=E(a_1X_1+a_2X_2)Y=a_1\langle X_1,Y\rangle+a_2\langle X_2,Y\rangle.$
\item Symmetry: $\langle X,Y\rangle=EXY=EYX=\langle Y,X\rangle.$
\item PSD: $\langle X,X\rangle=EX^2\ge 0.$
\end{enumerate}
\end{lemma}

\subsection{Markov and Chebyshev Inequalities}
\begin{theorem}{Markov's Inequality}
Let $X$ be a random variable, let $X\ge 0$ and $t\ge 0$, then for nondecreasing $g:\mathbb{R}_{\ge 0}\to\mathbb{R}_{\ge 0}$, we have:
\[
P(X\ge t)\le\frac{Eg(X)}{g(t)}.
\]
\end{theorem}

\begin{proof}
Take an arbitrary nonnegative random variable $X\ge 0$, an arbitrary $t\in\mathbb{R}_{\ge 0}$ and a nondecreasing function $g:\mathbb{R}_{\ge 0}\to\mathbb{R}_{\ge 0}$, then obviously:
\[
g(t)I(X\ge t)\le g(X).
\]
Apply $E$ to both sides:
\[
g(t)P(X\ge t)\le Eg(X)
\]
Hence, $\displaystyle P(X\ge t)\le\frac{Eg(X)}{g(t)}$.
\end{proof}
\newpage
\begin{corollary}[Chebyshev Inequalities]
Let $X\in L_2$, $\varepsilon>0$, $\mu=EX$, $\sigma=\sqrt{(X-EX)^2}$, and $\displaystyle k=\frac{\varepsilon}{\sigma}>0$, then we have:
\begin{enumerate}
\item Chebyshev I: 
\[
\displaystyle P(|X-\mu|\ge \varepsilon)\le\frac{\sigma^2}{\varepsilon^2}.
\]
\item Chebyshev II: 
\[
\displaystyle P\left(\frac{X-\mu}{\sigma}\ge k\right)\le \frac{1}{k^2}.
\]
\end{enumerate}
\end{corollary}

\begin{proof}[Proof of Chebyshev I]
Let $X\in :_2$ and $\mu=EX$. Let $Y=(X-\mu)^2$, then $EY=\sqrt{(X-EX)^2}=\sigma^2$. Hence:
\[
P(|X-\mu|\ge \varepsilon)=P(|X-\mu|^2\ge\varepsilon^2)\le\frac{E(x-\mu)^2}{\varepsilon^2},
\]
since $\sigma=\sqrt{(X-EX)^2}$:
\[
P(|X-\mu|\ge\varepsilon)\le\frac{E(x-\mu)^2}{\varepsilon^2}=\frac{\sigma^2}{\varepsilon^2}.
\]
\end{proof}

\begin{proof}[Proof of Chebyshev II]
Let $\displaystyle k=\frac{\varepsilon}{\sigma}$, then:
\[
P\left(\frac{X-\mu}{\sigma}\ge k\right)=P\left(\frac{X-\mu}{\sigma}\ge\frac{\varepsilon}{\sigma}\right)=P(X-\mu\ge\varepsilon)
\]
By Chebyshev I:
\[
P\left(\frac{X-\mu}{\sigma}\right)=P(X-\mu\ge\varepsilon)\le P(|X-\mu|\ge\varepsilon)\le\frac{\sigma^2}{\varepsilon^2}=\frac{1}{k^2}.
\]
\end{proof}
\begin{definition}{$X$ Equals $Y$ wP1}
Let $X, Y$ be two random variables.
We say $X \overset{\text{wP1}}{=} Y$ iff
\[
P(\{\omega \in \Omega : X(\omega) = Y(\omega)\}) = 1,
\]
i.e.\ $P(X = Y) = 1$.
\end{definition}

\begin{corollary}[Markov's Equality]
$X \overset{\text{wP1}}{=} 0\iff E|X|=0$
\end{corollary}

\subsection{The Geometry of $L_2$}
\begin{definition}{Orthogonality}
Let $X,Y\in L_2$, then we say $X$ is orthogonal to $Y$, denoted $X\perp Y$, if: 
\[
\|Y-X\|=\|Y+X\|, \text{ i.e. }\langle X,Y\rangle=EXY=0.
\]
\end{definition}

\begin{definition}{Orthogonal Projection}
Let $Y\in L_2$, then $\widehat Y$ is the orthogonal projection of $Y$ on $X$ if:
\begin{enumerate}
\item $\widehat Y=tx$ for some $t\in\mathbb{R}$.
\item $Y-\widehat Y\perp X$.
\end{enumerate}
\end{definition}

\begin{definition}{The Cosine of the Angle Between Two Random Variables}
Let $\theta$ be the angle between two random variables $X,Y\in L_2$, then:
\[
\cos\theta(X,Y)=\frac{EXY}{\sqrt{EX^2EY^2}}.
\]
\end{definition}
\begin{lemma}[Cauchy-Schwarz Inequalities]
There are three variations of the Cauchy-Schwarz Inequality:
\begin{enumerate}
\item Cauchy-Schwarz I: $\|\widehat Y\|\le\|Y\|\iff Y\overset{\text{wP1}}{=}\widehat Y $.
\item Cauchy-Schwarz II: $EXY\le\sqrt{EX^2EY^2} \iff Y\overset{\text{wP1}}{=}tX\text{ for some }t\in\mathbb{R}$.
\item Cauchy-Schwarz III: $|\cos\theta(x,y)|\le 1\iff Y\overset{\text{wP1}}{=}tX\text{ for some }t\in \mathbb{R}$.
\end{enumerate}
\end{lemma}

\begin{lemma}[Triangle Inequality]
$\|X+Y\|\le\|X\|+\|Y\| \iff Y\overset{\text{wP1}}{=}tX\text{ for some }t\in \mathbb{R_{\ge 0}}$.
\end{lemma}

\subsection{Mean, Variance and Standard Deviation}

\subsection{Covariance and Correlation}

\subsection{The Linear Model}

\newpage
% =====================================================================
\section{Types of Distributions}
% =====================================================================

\subsection{Joint Distribution and Statistical Independence}
\begin{definition}{Statistical Independence}
Let $X,Y$ be random variables, then $X$ is statistical independent to $Y$, denoted $X\perp\!\!\!\!\perp Y$ iff:
\[
P(X\in A|Y\in B)=P(X\in A).
\]
\end{definition}

\begin{proposition}[Properties of Statistical Independence]

\end{proposition}

\subsection{Gamma Distribution}

\subsection{Beta Distribution}

\subsection{Binomial Distribution}

\subsection{Multinomial Distribution}

\subsection{Geometric and Hypergeometric Distribution}

\subsection{Normal Distribution}

\newpage
% =====================================================================
\section{Types of Convergence}
% =====================================================================

\subsection{L-Convergence}

\subsection{Convergence in Probability}

\subsection{Convergence with Probability 1}

\subsection{Convergence in Distribution}

\subsection{Weak Convergence}

\newpage
% =====================================================================
\section*{Notes \& Questions}
\addcontentsline{toc}{section}{Notes \& Questions}
% =====================================================================

\begin{tcolorbox}[colback=black!2, colframe=neutralColor!40, arc=1.5mm, breakable]
\begin{itemize}[leftmargin=*]
  \item[$\square$] 
  \item[$\square$] 
  \item[$\square$] 
\end{itemize}
\end{tcolorbox}

\end{document}

